% French original: extra/SGA-archive/sga7-1.pdf, physical PDF pages 2--4.
% This native translation supplements, and does not alter, the accepted187 reader components.
% Original introduction first page is unnumbered; following pages bear VI and VII.
\phantomsection
\pdfbookmark[0]{导言}{sga7i-introduction}
\thispagestyle{empty}
\begin{center}
\underline{导言}
\end{center}

% FM-P2.1
设 $S$ 是概形，$f:X\longrightarrow S$ 是概形态射。如果 $f$ 是固有光滑态射，且素数 $\ell$ 在 $S$ 上可逆，那么 $f$ 的几何纤维的 $\ell$ 进上同调群 $H^i(X_{\bar s},\mathbb Z_\ell)$ 在 $S$ 上构成一个 $\ell$ 进局部系统。如果只假定 $f$ 固有，这些群就是 $S$ 上一个 $\ell$ 进层 $R^if_*\mathbb Z_\ell$ 的纤维。\underline{消灭圈理论}把这个层在 $S$ 上的分歧与 $f$ 的奇点联系起来。

% FM-P2.2
我们只考虑 $S$ 为 Hensel 性状的情形，即 $S$ 是一个 Hensel 离散赋值环的谱。在实际应用中，这就是主要情形。关于这一理论的启发式描述，参见 (I~2.1)。设 $S=\Spec(V)$，$k(\bar\eta)$ 是 $V$ 的分式域 $k(\eta)$ 的一个代数闭包，$k(\bar s)$ 是剩余域 $k(s)$ 的相应代数闭包；这里 $k(\bar s)$ 是 $V$ 在 $k(\bar\eta)$ 中的整闭包 $\bar V$ 的剩余域。在 $X$ 在 $S$ 上固有、平坦、相对维数为 $n$，并且 $f$ 只有一个非光滑点 $x\in X_s$ 的特殊情形中，我们定义 $\Gal(\bar\eta/\eta)$ 模 $\phi^i$。这些模纯粹由 $x$ 的邻域中的局部情形决定，当 $i\notin[0,n]$ 时为零 (I~4.2)。我们还构造一个 $\Gal(\bar\eta/\eta)$ 模的长正合列
\[
\cdots\longrightarrow H^i(X_{\bar s},\mathbb Z_\ell)
\xrightarrow{\mathrm{sp}}H^i(X_{\bar\eta},\mathbb Z_\ell)
\longrightarrow\phi^i\longrightarrow
H^{i+1}(X_{\bar s},\mathbb Z_\ell)\longrightarrow\cdots.
\]
这里把 $\Gal(\bar s/s)$ 模 $H^i(X_{\bar s},\mathbb Z_\ell)$ 看作惯性群作用平凡的 $\Gal(\bar\eta/\eta)$ 模；$\mathrm{sp}$ 是专门化映射。

% FM-P2.3
我们还给出使低次的 $\phi^i$ 为零的判据，但这些判据目前仅在特征 $0$ 下有效。例如，如果再假设 $X$ 是局部完全交，那么当 $i\ne n$ 时，$\phi^i=0$ (I~4.5)。

\clearpage
% FM-P3.1
\underline{单值定理}断言，$\Gal(\bar\eta/\eta)$ 的惯性子群 $I$ 的作用是拟幺幂的：对每个 $T\in I$，存在整数 $N>0$、$M>0$，使 $H^i(X_{\bar\eta},\mathbb Z_\ell)$ 的自同态 $(T^M-1)^N$ 为零。我们给出这个定理的两个证明。第一个证明 (I~1.2) 具有算术性质，只要所考虑的上同调群满足适当的有限性性质，就可以应用。关键在于：当 $k(s)$ 在其素域上有限生成时，对于 $\Gal(\bar\eta/\eta)$ 的任意 $\ell$ 进表示，$I$ 的作用都是拟幺幂的。第二个证明更具几何性质，需要纯性和奇点解消；目前，它只在特征 $0$ 下或者对 $H^1$ 有效。它给出了关于幂零指数 $N$ 的宝贵信息：对于 $H^i$，有 $N\leq i+1$。正如以后将看到的，在 $\mathbb C$ 上，这个证明很适合与超越理论进行比较。

% FM-P3.2
把这些结果应用于阿贝尔簇的 $H^1$，也就是它们的 Tate 模，就可以研究阿贝尔簇的模 $p$ 约化。我们给出阿贝尔簇稳定约化定理的两个证明。这个定理说，经过分歧扩张后，N\'eron 模型的连通特殊纤维是一个以环面为核、以阿贝尔簇为商的扩张。第一个证明 (I~6) 沿用单值定理的算术证明方法。第二个证明 (IX~3.6) 则依靠对 N\'eron 模型及其极化性质细致得多的分析。

% FM-P3.3
Grothendieck 的报告 I 至 V 没有写成文稿。它们已被概括成报告 I。报告 I 对所陈述结果的证明是简洁的，但基本上是完整的。

% FM-P3.4
报告 II 把 Lefschetz 铅笔的方法 (I~5.3) 应用于基本群的研究。对于代数闭域 $k$ 上的曲面 $S$，设 $k$ 的特征指数为 $p$；我们证明：基本群的素于 $p$ 完备化，即亲有限群 $\pi_1^{(p)}(S,s)$，在素于 $p$ 的亲有限群范畴中有限表示。

% FM-P3.5
如上所述，报告 III 至 V 并不存在。

\clearpage
% FM-P4.1
报告 VI 叙述 Schlessinger 的形变理论，并作了一些补充。其中特别注意了所分类对象的无穷小自同构，避免过于仓促地转到对象的同构类函子。报告还给出了 $\underline{\Ext}^{i}(L_{X/S},-)$ 的一种新构造，其中 $L_{X/S}$ 是 $X/S$ 的相对余切复形。这个报告在本讨论班后续部分的用途，主要在于其中对普通二次奇点形变的研究。

% FM-P4.2
报告 VII 和 VIII 讨论双扩张理论。这个理论在报告 IX 中起着关键作用，用来表述从阿贝尔簇转到它的 N\'eron 模型时，极化发生了什么变化。报告 IX 包含阿贝尔簇稳定约化定理的证明，以及各种应用。

% FM-P4.3
本讨论班的续卷 SGA~7~II 由 P.~Deligne 和 N.~Katz 编写，将于以后出版。

\vspace{1.5em}
\begin{flushright}
1972年5月，Bures sur Yvette\\[0.5em]
P. DELIGNE
\end{flushright}
